\chapter{IIoT, OPC UA e MQTT}
\label{cap:iiot}

\begin{objetivos}
Ao final deste capítulo você deve ser capaz de:
\begin{itemize}
  \item explicar o que a arquitetura cliente-servidor clássica não resolve;
  \item descrever o que OPC UA acrescenta a um protocolo de transporte;
  \item explicar o modelo \emph{publish/subscribe} do MQTT e o papel dos níveis de serviço;
  \item dizer o que o Sparkplug padroniza e por que isso importa;
  \item apontar os limites do IIoT: latência, determinismo, custo de rede e borda versus nuvem.
\end{itemize}
\end{objetivos}

\section{O que a automação clássica deixou em aberto}

O Capítulo~\ref{cap:estrutura} mostrou a hierarquia em que cada nível fala com o de cima por um
caminho próprio. Esse arranjo funcionou por décadas e tem duas limitações que ficaram caras com o
tempo:

\begin{enumerate}
  \item \textbf{o problema $N \times M$}: cada combinação de aplicação e equipamento exige um
        \emph{driver} próprio. Três sistemas de supervisão e cinco famílias de controladores
        significam até quinze integrações para manter, atualizar e testar a cada versão de
        \emph{firmware};
  \item \textbf{dado sem significado}: um protocolo de transporte diz \emph{onde} está o valor, não
        \emph{o que} ele é. Sem essa camada, cada consumidor precisa de uma tabela mantida à mão
        dizendo que o registro 40001 é a temperatura e que a unidade é grau Celsius.
\end{enumerate}

O conjunto de tecnologias que ataca esses dois problemas ficou conhecido como \textbf{IIoT}
(\emph{Industrial Internet of Things}): usar padrões de comunicação abertos para levar dado
industrial a qualquer consumidor autorizado, dentro ou fora da fábrica, sem que cada consumidor
precise conhecer o equipamento que o originou.

\index{IIoT}\index{Internet das Coisas!industrial}

\section{OPC UA: camada de informação, não apenas de transporte}

OPC UA (\emph{Unified Architecture}) foi lançado em 2008 pela OPC Foundation como uma arquitetura
orientada a serviços, independente de plataforma, que reuniu as especificações anteriores do OPC
``clássico'' em um único conjunto. Quatro características o distinguem dos protocolos de transporte
apresentados no capítulo anterior (OPC Foundation, \emph{Unified Architecture}, especificação
disponível publicamente):

\begin{itemize}
  \item \textbf{modelagem de informação}: o dado é exposto como um espaço de endereços hierárquico,
        com objetos, variáveis e tipos definidos. Um cliente pode \emph{descobrir} a estrutura, em
        vez de depender de uma tabela externa;
  \item \textbf{assinatura de dados}: o cliente se inscreve em um item e recebe notificação quando o
        valor muda (relatório por exceção), em vez de perguntar em laço;
  \item \textbf{eventos e métodos}: além de valores, o servidor notifica eventos e pode executar
        métodos, ou seja, há canal para invocar ações e não só para ler e escrever;
  \item \textbf{sessão segura}: autenticação por certificado, cifragem das mensagens, assinatura e
        sequenciamento, com registro de auditoria.
\end{itemize}

O detalhe que interessa a quem compra equipamento é o primeiro item da lista acima. Sobre a base do
OPC UA existe um programa de perfis setoriais (\emph{companion specifications}), mantido por
organizações de cada segmento industrial. Um perfil publicado define, para um tipo de equipamento,
\emph{quais} variáveis devem existir e com que nome --- que é exatamente o que faltava no Modbus, como
observado no Capítulo~\ref{cap:familia-iec}.

O mesmo padrão também prevê uma variante de publicação e assinatura (PubSub), pensada para
comunicação de muitos para muitos, em que o produtor publica sem saber quem consome. Guarde essa
distinção: é a mesma ideia que o MQTT explora, por outro caminho.

\index{OPC UA}

\section{MQTT: publicação e assinatura para telemetria}

MQTT é um padrão aberto de mensagens mantido pela OASIS (versões 3.1.1 e 5.0), projetado para
dispositivos com pouca memória e rede de banda limitada ou instável. O desenho é
\emph{publish/subscribe}:

\begin{itemize}
  \item um \textbf{broker} central recebe as mensagens publicadas;
  \item quem publica envia para um \textbf{tópico}, sem saber quem recebe;
  \item quem consome assina o tópico que lhe interessa, sem saber quem publica.
\end{itemize}

Esse desacoplamento é a razão do sucesso do MQTT em telemetria. Acrescentar um consumidor novo ---
um painel de indicadores, uma ferramenta de análise, um sistema de manutenção --- não exige tocar no
publicador. Em uma arquitetura ponto a ponto, cada consumidor novo é uma integração nova.

Dois recursos merecem atenção:

\begin{itemize}
  \item \textbf{níveis de serviço} (\emph{QoS} 0, 1 e 2, na documentação da OASIS): entregar no máximo
        uma vez, ao menos uma vez e exatamente uma vez. Escolher o nível é uma decisão de engenharia:
        contagem de peças tolera repetição; consumo acumulado de energia, não;
  \item \textbf{sessão persistente}: o cliente pode reconectar e receber o que foi publicado durante
        a ausência. É o que viabiliza o padrão \emph{store and forward} em enlaces que caem, comuns em
        rede celular e em campo remoto.
\end{itemize}

\section{Sparkplug: padronizar o conteúdo, não só o transporte}

MQTT padroniza \emph{como} a mensagem viaja, não \emph{o que} ela diz. Sem essa camada, cada
integração inventa seu formato de tópico e seu formato de carga útil, e volta-se ao problema
$N \times M$ --- agora em cima de um protocolo aberto.

O \textbf{Eclipse Sparkplug} (versão 3.0, mantida pela Eclipse Foundation) resolve isso definindo,
para o mundo industrial, a estrutura dos tópicos e o formato da carga útil, além de dois mecanismos
que valem a pena conhecer:

\begin{itemize}
  \item \textbf{certificados de nascimento e de morte}: ao conectar-se, o dispositivo publica seu
        estado e a lista do que vai informar; ao cair, é possível detectar a ausência. O consumidor
        sabe \emph{se} o dado é válido, não apenas se ele chegou;
  \item \textbf{consciência de estado} (\emph{state awareness}): a informação de qualidade acompanha o
        dado, o que permite distinguir ``valor parado'' de ``equipamento desligado''.
\end{itemize}

\section{A pilha completa}

\begin{figure}[H]
  \centering
  \includegraphics[width=0.98\linewidth]{18-pilha-iiot.pdf}
  \caption{Pilha típica de IIoT: controlador, borda, broker e aplicações.}
  \label{fig:pilha-iiot}
\end{figure}

A Figura~\ref{fig:pilha-iiot} mostra a montagem mais comum:

\begin{enumerate}
  \item o \textbf{controlador} expõe dados por OPC UA. Em instalações antigas, o caminho é um
        \emph{driver} do protocolo proprietário dentro do gateway;
  \item o \textbf{gateway de borda} é o tradutor: fala OPC UA (ou Modbus) de um lado e MQTT do outro,
        mantém um \emph{buffer} local e aplica alguma regra de filtro --- mudança de valor, limite de
        taxa, agregação;
  \item o \textbf{broker} distribui;
  \item as \textbf{aplicações} consomem: histórico, painéis, análise, manutenção e, quando fizer
        sentido, o próprio MES.
\end{enumerate}

Note o que \emph{não} está na figura: nenhum caminho de comando descendo pelo broker até o laço de
controle. Essa ausência é deliberada e é o conteúdo da próxima seção.

\index{gateway de borda}\index{broker MQTT}

\section{Vantagens do IIoT na automação}

\begin{table}[H]
  \centering
  \caption{O que o IIoT entrega, em termos práticos.}
  \label{tab:iiot-vantagens}
  \begin{tabular}{p{4.4cm} p{9.9cm}}
    \toprule
    Vantagem & Efeito prático \\
    \midrule
    Acesso a mais variáveis & Parâmetros de diagnóstico que já existiam no equipamento e não eram
      lidos passam a ser visíveis --- sem alterar a lógica de controle \\
    Integração desacoplada & Um consumidor novo entra por assinatura, sem alterar o publicador nem
      o programa do controlador \\
    Dado com contexto & Modelo de informação e carga útil padronizada reduzem a dependência de
      tabelas de registradores mantidas à mão \\
    Efeito de escala & Muitas máquinas podem ser acompanhadas por um número pequeno de
      especialistas, em vez de exigir presença física em cada ponto \\
    Histórico para a melhoria & Tendência, consumo e contagem de ciclos alimentam manutenção e
      eficiência energética, temas dos Capítulos~\ref{cap:dados} e~\ref{cap:eficiencia} \\
    \bottomrule
  \end{tabular}
\end{table}

\section{Desafios e limites}

Se o IIoT fosse apenas vantagem, toda planta já o teria adotado por completo. Não é o caso. Os
obstáculos reais são de quatro naturezas:

\begin{enumerate}
  \item \textbf{latência e determinismo}. A cadeia inclui gateway, broker, fila, rede compartilhada e
        aplicação. Cada elo acrescenta atraso e variabilidade. Um caminho que passa por broker e por
        rede corporativa não é determinístico --- e, por isso, não serve para controle. O critério
        prático: se a função precisa de prazo garantido, ela fica no controlador;
  \item \textbf{banda e custo de rede}. Enviar toda variável a cada segundo, de cada máquina, custa
        tráfego, licença e nuvem. O gateway existe também para isso: filtrar, agregar e transmitir por
        exceção. Sem essa disciplina, o projeto de IIoT entrega uma conta mensal e não um indicador;
  \item \textbf{segurança}. Cada caminho novo de leitura é superfície nova de ataque. Um broker
        acessível a partir da rede corporativa, com credencial padrão e sem TLS, transforma a melhoria
        em porta de entrada (Capítulo~\ref{cap:seguranca});
  \item \textbf{governança do dado}. Quem é o dono da variável? Quem aprova mudar o seu significado?
        Duas áreas publicando a mesma grandeza com nomes e unidades diferentes produzem dois
        indicadores e nenhuma decisão.
\end{enumerate}

\begin{atencao}
O IIoT amplia a \emph{visão} sobre o processo; não amplia a \emph{autoridade} sobre a máquina. Nada
que dependa de prazo garantido --- intertravamento, parada de emergência, proteção de máquina ---
pode depender de broker, de gateway ou de nuvem. A fronteira correta é simples de enunciar: dado
sobe, comando de segurança não desce por esse caminho. O Capítulo~\ref{cap:ti-to} retoma essa
fronteira ao tratar de tecnologias de TI no chão de fábrica.
\end{atencao}

\section{Por onde começar}

A ordem que costuma dar certo em planta em operação:

\begin{enumerate}
  \item \textbf{inventariar} o que já existe: quais controladores, quais protocolos, quais dados já
        são coletados e quem usa cada um;
  \item começar por um \textbf{indicador que a planta quer} --- consumo por linha, disponibilidade de
        um equipamento crítico --- e trabalhar de trás para frente, para saber quais variáveis
        faltam;
  \item \textbf{piloto em um equipamento}, com o gateway isolado da rede de controle;
  \item só então padronizar, documentando tópicos, nomes e unidades.
\end{enumerate}

Começar pela tecnologia --- ``vamos instalar a plataforma'' --- produz projeto bonito e nenhum
indicador. Começar pelo indicador produz uma lista de variáveis, que é o que a engenharia precisa.

\begin{pratica}
Escolha um equipamento da sua planta ou da sua bancada e responda, por escrito:
(a) quais variáveis de processo ele já disponibiliza em rede hoje?
(b) quais delas interessariam a um indicador de disponibilidade?
(c) qual seria o caminho do dado, no formato da Figura~\ref{fig:pilha-iiot}, até esse indicador?
(d) qual dessas variáveis \emph{não} pode ser levada para fora da rede de controle, e por quê?
\end{pratica}

\begin{exercicios}
\begin{enumerate}[label=\alph*)]
  \item Explique o problema $N \times M$ com um exemplo numérico de três sistemas e quatro famílias
        de equipamentos.
  \item Qual das quatro características do OPC UA resolve diretamente o problema do significado do
        dado? Justifique.
  \item Compare OPC UA e MQTT quanto a: modelo de comunicação, papel do intermediário e adequação a
        controle de tempo real.
  \item Descreva o caminho de um dado de temperatura desde o CLP até um painel de indicadores na
        nuvem, nomeando cada elemento da Figura~\ref{fig:pilha-iiot}.
  \item Por que a escolha do nível de serviço (QoS) muda conforme a variável? Dê um exemplo de
        variável que tolera repetição e um que não tolera.
  \item O que os certificados de nascimento e de morte informam ao consumidor que um simples valor
        recebido não informa?
  \item Cite dois riscos concretos de ligar o broker MQTT diretamente na rede corporativa.
  \item Um fornecedor propõe fazer o intertravamento de uma prensa ``pelo servidor OPC UA, para
        centralizar a lógica''. Escreva a recusa técnica em cinco linhas.
  \item Escolha uma das vantagens da Tabela~\ref{tab:iiot-vantagens} e diga qual desafio da seção
        seguinte pode anulá-la se o projeto for mal feito.
  \item Descreva, em quatro passos, como você conduziria um piloto de IIoT em uma máquina isolada.
\end{enumerate}
\end{exercicios}
